\chapter{Designs and Finite Geometries}

\begin{center}\itshape
Fifteen young ladies in a school walk out three abreast for seven days in succession: it is required to arrange them daily, so that no two shall walk twice abreast.\\
\upshape(Thomas Kirkman, \emph{The Lady's and Gentleman's Diary}, 1850)
\end{center}

The Reverend Thomas Kirkman was the rector of Croft, a village in Lancashire, for more than fifty years, and he did
mathematics in the time his parish left him. In 1847 he answered a question from the same magazine: for which $n$ can one
choose triples from $n$ objects so that every pair of objects lies in exactly one triple? He proved that this is possible
exactly when $n$ leaves remainder $1$ or $3$ on division by $6$. Six years later Jakob Steiner, in Berlin, asked the same question
without knowing Kirkman's paper, and the objects are called \emph{Steiner} triple systems to this day. Kirkman's schoolgirls are a
harder version: the $35$ triples of fifteen girls must also be split into seven days, each day using every girl once. A solution is
in Table~\ref{tab:kirkman}; the general problem, for $n\equiv3\pmod 6$ schoolgirls, waited until 1971.

Such arrangements are called \emph{designs}, and their most important users were not schoolmistresses but farmers. At the Rothamsted
agricultural station in the 1920s and 1930s, Ronald Fisher and then Frank Yates planned field trials in which varieties of a crop were compared on plots
grouped into blocks, and to compare fairly they needed every pair of varieties to share the same number of blocks. In Calcutta, R.~C.~Bose
found that the best such plans come from geometry over finite fields. The same objects are now used to test software: if a program has
twenty settings and bugs are usually caused by one or two of them interacting, a well-designed set of a few dozen tests covers every pair
of values of every pair of settings.

This chapter builds these objects. Steiner triple systems exist for all the allowed $n$, by two explicit constructions. Finite fields give
projective planes, in which any two points lie on one line and any two lines meet in one point, and Latin squares that are pairwise
orthogonal; the two are the same thing in different clothes. Euler believed that there are no two orthogonal Latin squares of order $6$,
and was right, but his reason was wrong; a theorem of Bruck and Ryser, using Lagrange's four squares, shows that there is no projective
plane of order $6$ or $14$, and nobody knows whether one of order $12$ exists. We end with the covering arrays of software testing.

\begin{think}
Sixteen cards: the aces, kings, queens and jacks of the four suits. Arrange them in a $4\times4$ square so that every row and every
column contains one card of each rank and one card of each suit. Then try the same with the nine cards of ranks $A,K,Q$ in three suits,
and with $36$ officers of six ranks from six regiments.
\end{think}

\section{Designs}

\begin{definition}\label{def:design}
Let $v>k\ge2$ and $\lambda\ge1$. A \emph{$2$-$(v,k,\lambda)$ design} is a set $X$ of $v$ \emph{points} with a family of $k$-element subsets, the \emph{blocks}, such
that every pair of distinct points is contained in exactly $\lambda$ blocks. A \emph{Steiner triple system} of order $v$, an $\mathrm{STS}(v)$, is a $2$-$(v,3,1)$ design.
\end{definition}

The Fano plane (Figure~\ref{fig:fano}) is an $\mathrm{STS}(7)$: its seven lines $124,135,236,167,257,347,456$ are blocks, and any two points lie on one line. The
twelve lines of $\mathbb F_3^2$ (Figure~\ref{fig:caps}) form an $\mathrm{STS}(9)$.

\begin{proposition}\label{prop:params}
In a $2$-$(v,k,\lambda)$ design, every point lies in the same number $r$ of blocks, and if $b$ is the number of blocks,
\[
r(k-1)=\lambda(v-1),\qquad bk=vr.
\]
\end{proposition}
\begin{proof}
Fix a point $x$ and count the pairs $(y,B)$ with $y\ne x$ and $\{x,y\}\subseteq B$. Each of the $v-1$ points $y$ gives $\lambda$ blocks; each of the $r_x$ blocks through $x$
gives $k-1$ points $y$. So $r_x(k-1)=\lambda(v-1)$, which does not depend on $x$. Counting the pairs $(x,B)$ with $x\in B$ gives $bk=vr$.
\end{proof}

So the parameters are not free: $k-1$ must divide $\lambda(v-1)$, and $k$ must divide $vr$. For triple systems, $r=\frac{v-1}2$ and $b=\frac{v(v-1)}6$.

\begin{theorem}[Fisher, 1940]\label{thm:fisherdesign}
In a $2$-$(v,k,\lambda)$ design, $b\ge v$.
\end{theorem}
\begin{proof}
For each point $x$, let $S_x$ be the set of blocks containing $x$, a subset of the set of $b$ blocks. For $x\ne y$, $|S_x\cap S_y|=\lambda\ge1$. The sets $S_x$ are distinct:
if $S_x=S_y$ then $r=|S_x|=|S_x\cap S_y|=\lambda$, and $r(k-1)=\lambda(v-1)$ gives $k=v$, which the definition excludes. By Fisher's inequality in the form of Theorem~\ref{thm:fisher},
there are at most $b$ such sets: $v\le b$.
\end{proof}

A design with $b=v$ is called \emph{symmetric}; then $r=k$, and any two blocks meet in exactly $\lambda$ points (Exercise~\ref{ex:symmetric}). The projective planes of
Section~\ref{sec:planes} are the symmetric designs with $\lambda=1$.

\section{Steiner triple systems}

\begin{proposition}\label{prop:stsnec}
If an $\mathrm{STS}(v)$ exists, then $v\equiv1$ or $3\pmod 6$.
\end{proposition}
\begin{proof}
By Proposition~\ref{prop:params}, $r=\frac{v-1}2$ and $b=\frac{v(v-1)}6$ are integers. So $v$ is odd, and $v\bmod6\in\{1,3,5\}$. If $v=6m+5$, then $v(v-1)=(6m+5)(6m+4)\equiv20\equiv2\pmod 6$,
and $b$ is not an integer.
\end{proof}

The converse needs constructions, and the two classical ones use a small algebraic gadget. A \emph{commutative quasigroup} of order $n$ is a set $Q$ of $n$ elements with an operation
$\circ$ such that $x\circ y=y\circ x$ and, for each $x$, the map $y\mapsto x\circ y$ is a bijection of $Q$; its table is a symmetric Latin square.

\begin{theorem}[Bose, 1939]\label{thm:bose}
For every $m\ge0$ there is an $\mathrm{STS}(6m+3)$.
\end{theorem}
\begin{proof}
Let $n=2m+1$ and $Q=\Z/n\Z$ with $x\circ y=(m+1)(x+y)$. Since $2(m+1)=n+1\equiv1$, this is ``$\frac{x+y}2$'': it is commutative, $y\mapsto x\circ y$ is a bijection, and $x\circ x=x$. Take as points
the $3n=6m+3$ pairs $(x,i)$ with $x\in Q$ and $i\in\Z/3\Z$ (three copies of $Q$, the \emph{layers}), and as blocks
\begin{enumerate}[label=(\Alph*)]
\item $\{(x,0),(x,1),(x,2)\}$ for $x\in Q$;
\item $\{(x,i),(y,i),(x\circ y,i+1)\}$ for $x\ne y$ in $Q$ and $i\in\Z/3\Z$.
\end{enumerate}
A pair of points in the same layer, $\{(x,i),(y,i)\}$, lies in exactly one block, of type (B) with layer $i$; blocks of type (A) and the other blocks of type (B) contain no
pair inside layer $i$. Any other pair is $\{(a,i),(b,i+1)\}$ for some $i$, and the blocks containing it are: the block (A) if $a=b$, and the blocks (B) of layer $i$ with
$x=a$ and $a\circ y=b$, $y\ne a$. There is exactly one $y$ with $a\circ y=b$, and it differs from $a$ exactly when $b\ne a\circ a=a$. So the pair lies in exactly one block.
\end{proof}

For $v=6m+1$ the idea is the same, with one extra point $\infty$ and an even number of elements in the quasigroup, which then cannot be idempotent (Exercise~\ref{ex:idempotent}).

\begin{lemma}\label{lem:halfidem}
Let $n\ge1$ and $Q=\Z/2n\Z$. Let $\sigma:Q\to Q$ send $2j$ to $j$ and $2j+1$ to $n+j$, for $0\le j<n$, and let $x\circ y=\sigma(x+y)$. Then $(Q,\circ)$ is a commutative quasigroup
with $x\circ x=(n+x)\circ(n+x)=x$ for $0\le x<n$.
\end{lemma}
\begin{proof}
$\sigma$ is a bijection (it sends the even residues onto $\{0,\dots,n-1\}$ and the odd ones onto $\{n,\dots,2n-1\}$), and $y\mapsto x+y$ is a bijection, so $y\mapsto x\circ y$ is one.
Commutativity is clear. For $0\le x<n$, $x+x=2x$ and $(n+x)+(n+x)\equiv2x$, and $\sigma(2x)=x$.
\end{proof}

\begin{theorem}[Skolem, 1958]\label{thm:skolem}
For every $m\ge0$ there is an $\mathrm{STS}(6m+1)$.
\end{theorem}
\begin{proof}
For $m=0$ there is nothing to do. For $m\ge1$, let $Q=\Z/2m\Z$ with the operation of Lemma~\ref{lem:halfidem} ($n=m$), and take as points $\infty$ and the $6m$ pairs $(x,i)$, $x\in Q$, $i\in\Z/3\Z$.
The blocks are
\begin{enumerate}[label=(\Alph*)]
\item $\{(x,0),(x,1),(x,2)\}$ for $0\le x<m$;
\item $\{\infty,(m+x,i),(x,i+1)\}$ for $0\le x<m$ and $i\in\Z/3\Z$;
\item $\{(x,i),(y,i),(x\circ y,i+1)\}$ for $x\ne y$ in $Q$ and $i\in\Z/3\Z$.
\end{enumerate}
\emph{Pairs with $\infty$.} The point $(a,i)$ is in the block (B) with $x=a-m$ and layer $i$ if $a\ge m$, and in the block (B) with $x=a$ and layer $i-1$ if $a<m$; in exactly one block in both cases.

\emph{Pairs in one layer.} As for Bose's construction, $\{(x,i),(y,i)\}$ lies only in the block (C) of layer $i$ with these $x,y$.

\emph{Pairs $\{(a,i),(b,i+1)\}$.} The candidates are: (A) when $a=b<m$; (B) when $a\ge m$ and $b=a-m$; (C) with $x=a$ and $y$ the unique element with $a\circ y=b$, valid when $y\ne a$, that is
when $b\ne a\circ a$. By Lemma~\ref{lem:halfidem}, $a\circ a=a$ if $a<m$ and $a\circ a=a-m$ if $a\ge m$. So if $a<m$: when $b=a$ only (A) applies, otherwise only (C). If $a\ge m$: when $b=a-m$ only (B)
applies, otherwise only (C). In every case the pair is in exactly one block.
\end{proof}

With Proposition~\ref{prop:stsnec}, an $\mathrm{STS}(v)$ exists if and only if $v\equiv1,3\pmod 6$: Kirkman's theorem of 1847, with constructions that the notebook runs. As
a check, the numbers of blocks are $m+3m+3\binom{2m}2=m(6m+1)=\frac{v(v-1)}6$ and $n+3\binom n2=\frac{n(3n-1)}2=\frac{v(v-1)}6$.

There are many triple systems: one each of orders $7$ and $9$ up to relabelling (Exercise~\ref{ex:sts7}), two of order $13$, $80$ of order $15$, and $11\,084\,874\,829$ of order
$19$ (Kaski and \"Osterg\aa rd, 2004), after which the numbers grow like $v^{v^2/6}$. A triple system is \emph{resolvable} if its blocks can be split into \emph{parallel classes},
each partitioning the points; Kirkman's schoolgirls ask for a resolvable $\mathrm{STS}(15)$, and Table~\ref{tab:kirkman} shows one, found by the notebook among the lines of the
projective space over $\mathbb F_2$ of dimension $3$. Resolvable systems need $v\equiv3\pmod 6$ (Exercise~\ref{ex:resolvable}), and Ray-Chaudhuri and Wilson proved in 1971 that this
suffices.

\begin{table}[ht]
\centering
\begin{tabular}{ll}
Monday & ABC\ \ DHL\ \ EJO\ \ FKM\ \ GIN\\
Tuesday & ADE\ \ BHJ\ \ CMN\ \ FIO\ \ GKL\\
Wednesday & AFG\ \ BIK\ \ CLO\ \ DJN\ \ EHM\\
Thursday & AHI\ \ BLN\ \ CEF\ \ DKO\ \ GJM\\
Friday & AJK\ \ BMO\ \ CDG\ \ EIL\ \ FHN\\
Saturday & ALM\ \ BDF\ \ CIJ\ \ EKN\ \ GHO\\
Sunday & ANO\ \ BEG\ \ CHK\ \ DIM\ \ FJL
\end{tabular}
\caption{A solution of Kirkman's problem: fifteen girls $A,\dots,O$, seven days, no two girls twice in the same row. Each of the $105$ pairs appears exactly once.}
\label{tab:kirkman}
\end{table}

\section{Projective and affine planes}\label{sec:planes}

\begin{definition}\label{def:pp}
A \emph{projective plane} is a set of points with a family of subsets called lines such that (P1) two distinct points lie on exactly one line, (P2) two distinct lines meet in exactly
one point, and (P3) there are four points no three of which are on one line.
\end{definition}

Axiom (P3) only excludes degenerate cases, such as a single line with all points on it. The axioms are symmetric under exchanging points and lines (dually, (P3) gives four lines
no three through a point), so every statement about projective planes has a dual statement, proved by the dual proof.

\begin{proposition}\label{prop:order}
In a finite projective plane there is an integer $n\ge2$, the \emph{order}, such that every line has $n+1$ points, every point lies on $n+1$ lines, and there are $n^2+n+1$ points and
$n^2+n+1$ lines.
\end{proposition}
\begin{proof}
If $P$ is a point not on the line $\ell$, the lines through $P$ correspond bijectively to the points of $\ell$: each line through $P$ meets $\ell$ in one point (P2), and each point of $\ell$
is joined to $P$ by one line (P1). So the number of lines through $P$ is $|\ell|$.

Let $\ell,m$ be two lines. We find a point on neither. Take four points as in (P3); each line contains at most two of them. If $\ell\cup m$ misses one of them, we are done. Otherwise
$\ell$ contains exactly two, say $A,B$, and $m$ the other two, $C,D$. Let $X$ be the common point of the lines $AC$ and $BD$. If $X$ were on $\ell$, then $X=A$ (the lines $AC$ and $\ell$ meet only
at $A$, as $C\notin\ell$), and $A$ would be on $BD$: $A,B,D$ collinear, a contradiction. Symmetrically $X\notin m$. Then $|\ell|$ and $|m|$ are both the number of lines through $X$: all lines
have the same size, $n+1$ say. Every point $P$ is off some line, so it lies on $n+1$ lines.

The $n+1$ lines through a point $P$ contain every other point exactly once (P1), and $n$ other points each: there are $1+(n+1)n=n^2+n+1$ points. Counting incident pairs (point, line)
gives $(n^2+n+1)(n+1)=b(n+1)$, so there are $n^2+n+1$ lines. Finally, with $A,B,C,D$ as in (P3), the lines $CA$, $CB$, $CD$ are distinct, so $n+1\ge3$.
\end{proof}

\begin{theorem}\label{thm:pg}
Let $F$ be a finite field with $q$ elements. The $1$-dimensional subspaces of $F^3$ as points, and the $2$-dimensional subspaces as lines (a point lying on a line when it is contained in it), form
a projective plane $\mathrm{PG}(2,q)$ of order $q$.
\end{theorem}
\begin{proof}
(P1) Two distinct $1$-dimensional subspaces span a $2$-dimensional one, the only one containing both. (P2) Two distinct $2$-dimensional subspaces $U,W$ satisfy $U+W=F^3$, so by the dimension
formula $\dim(U\cap W)=2+2-3=1$. (P3) The subspaces spanned by $e_1$, $e_2$, $e_3$ and $e_1+e_2+e_3$: any three of these vectors are independent. A $2$-dimensional subspace has $q^2-1$ non-zero vectors,
and each $1$-dimensional subspace in it has $q-1$ of them, so a line has $q+1$ points.
\end{proof}

Since a field with $q$ elements exists for every prime power $q$ (Tome~2, Theorem~9.6), there are projective planes of every prime power order. Concretely, a point of $\mathrm{PG}(2,q)$ is a
vector $(x:y:z)\ne0$ up to a non-zero factor, and a line is the set of points with $ax+by+cz=0$ for a fixed $(a:b:c)$: the incidence is $ax+by+cz=0$, symmetric in points and lines.
The Fano plane is $\mathrm{PG}(2,2)$; Chapter~3 counted its seven points with a $q$-binomial coefficient, and Chapter~8 used $\mathrm{PG}(2,p)$ to build graphs with many edges and no square.

Planes of prime power order not isomorphic to $\mathrm{PG}(2,q)$ exist, for instance of order $9$; but every known projective plane has prime power order. Whether this must be so is
one of the oldest open problems of combinatorics. Removing a line turns a projective plane into the flat plane of school geometry, with parallels.

\begin{definition}\label{def:affine}
An \emph{affine plane} is a set of points with a family of subsets called lines such that (A1) two distinct points lie on exactly one line, (A2) for every line $\ell$ and point $P\notin\ell$ there is exactly
one line through $P$ disjoint from $\ell$, and (A3) there are three points not on one line.
\end{definition}

\begin{proposition}\label{prop:affine}
\begin{enumerate}[label=(\alph*)]
\item Removing from a projective plane of order $n$ one line $\ell_\infty$ and its points gives an affine plane with $n^2$ points, $n$ points on each line, in which the lines fall into $n+1$
\emph{parallel classes} of $n$ pairwise disjoint lines, each class partitioning the points.
\item Conversely, every finite affine plane arises in this way from a projective plane.
\end{enumerate}
\end{proposition}
\begin{proof}
(a) The remaining lines lose one point each, so have $n$ points, and $n^2+n+1-(n+1)=n^2$ points remain. (A1) holds as before. Two remaining lines are disjoint exactly when they met on $\ell_\infty$.
Given $\ell$ meeting $\ell_\infty$ at $Q$ and $P\notin\ell$, the line $PQ$ is the unique line through $P$ disjoint from $\ell$: (A2). The $n$ lines other than $\ell_\infty$ through a point $Q\in\ell_\infty$ form a
parallel class, and each affine point lies on exactly one of them. (A3) follows from (P3) applied to points off $\ell_\infty$ (Exercise~\ref{ex:affinea3}).

(b) Call two lines \emph{parallel} if they are equal or disjoint. This is an equivalence relation: if $\ell\parallel m$, $m\parallel k$ and $\ell\ne k$ met at a point $P$, then $P\notin m$ and $\ell$, $k$ would be two
lines through $P$ disjoint from $m$, against (A2). Add one new point for each parallel class, lying on all the lines of that class, and one new line $\ell_\infty$ consisting of the new points. (P1): two old
points are on one old line; an old point $P$ and the new point of a class lie on the unique line of the class through $P$ (it exists by (A2) or because $P$ is on a line of the class, and is unique since
the class is a set of disjoint lines); two new points lie only on $\ell_\infty$. (P2): two old lines meet at an old point, or are parallel and meet at the new point of their class, and not both;
an old line meets $\ell_\infty$ at the point of its class. (P3): let $A,B,C$ be old points not on a line, and $D$ the common point of the line through $C$ parallel to $AB$ and the line through $B$
parallel to $AC$ (these lines are in different classes, since $AB$ and $AC$ meet, so they meet). If $D$ were on $AB$, the parallel to $AB$ through $C$ would be $AB$ itself, and $C\in AB$; likewise $D\notin AC$; and if $B,C,D$
were collinear, the parallel to $AC$ through $B$ would be $BC$, which meets $AC$ at $C$. So no three of $A,B,C,D$ are collinear.
\end{proof}

\begin{think}
In the affine plane $\mathbb F_q^2$, what are the parallel classes? How many points does the plane have ``at infinity'', and what do they correspond to?
\end{think}

\section{Latin squares}

\begin{definition}\label{def:latin}
A \emph{Latin square} of order $n$ is an $n\times n$ array with entries in an $n$-element set of symbols, each symbol appearing once in each row and once in each column. Two Latin squares $L,M$
of order $n$ are \emph{orthogonal} if the $n^2$ pairs $(L(i,j),M(i,j))$ are all different, that is, each of the $n^2$ possible pairs appears exactly once. A family of Latin squares is
\emph{mutually orthogonal} (MOLS) if any two of them are orthogonal. We write $N(n)$ for the largest size of a family of MOLS of order $n$.
\end{definition}

The card puzzle of the opening, printed in the 1725 edition of Ozanam's \emph{R\'ecr\'eations math\'ematiques}, asks for two orthogonal Latin squares of order $4$, one for the ranks and one for the
suits; Figure~\ref{fig:cards} gives a solution. A Latin rectangle can always be completed to a Latin square (Chapter~6, Exercise~\ref{ex:latin}), so Latin squares are easy to find; orthogonal
ones are not.

\begin{figure}[ht]
\centering
\fig{ch12-cards}
\caption{Ozanam's puzzle: the sixteen court cards in a square, each rank and each suit once in every row and every column. Ranks and suits are two orthogonal Latin squares of order $4$, built
from the field $\mathbb F_4$ as in Corollary~\ref{cor:pmols}.}
\label{fig:cards}
\end{figure}

\begin{proposition}\label{prop:molsbound}
For $n\ge2$, $N(n)\le n-1$.
\end{proposition}
\begin{proof}
Renaming the symbols of a square keeps it Latin and keeps orthogonality. So in a family of MOLS we may assume that every square has first row $1,2,\dots,n$. Look at the entry in row $2$,
column $1$. It is not $1$, since $1$ is already in column $1$. Two squares cannot have the same entry $s$ there: the pair $(s,s)$ already appears in row $1$, column $s$. So the entries are distinct
elements of $\{2,\dots,n\}$.
\end{proof}

\begin{theorem}[Bose, 1938]\label{thm:planemols}
For $n\ge2$, there is a projective plane of order $n$ if and only if $N(n)=n-1$.
\end{theorem}
\begin{proof}
By Proposition~\ref{prop:affine}, we may work with affine planes of order $n$ ($n^2$ points, $n+1$ parallel classes of $n$ lines of $n$ points).

\emph{From a plane to squares.} Choose two parallel classes, the \emph{rows} $h_1,\dots,h_n$ and the \emph{columns} $c_1,\dots,c_n$. A row and a column are not parallel, so they meet in one
point, and every point is on one row and one column: the points are the pairs $(i,j)$, with $(i,j)=h_i\cap c_j$. For each of the other $n-1$ classes $\mathcal C$, number its lines $1,\dots,n$, and
let $L_{\mathcal C}(i,j)$ be the number of the line of $\mathcal C$ through $(i,j)$. Each line of $\mathcal C$ meets the row $h_i$ in exactly one point, so row $i$ of $L_{\mathcal C}$ contains every number
once; likewise for columns. For two classes $\mathcal C\ne\mathcal C'$, line $s$ of $\mathcal C$ and line $s'$ of $\mathcal C'$ meet in exactly one point, which is the unique cell where the pair $(s,s')$
appears. So the $n-1$ squares are mutually orthogonal.

\emph{From squares to a plane.} Let $L_1,\dots,L_{n-1}$ be MOLS of order $n$. Take the $n^2$ cells $(i,j)$ as points and as lines: the rows, the columns, and for each $k$ and each symbol $s$ the
set of cells where $L_k$ has the entry $s$. This gives $n+1$ classes of $n$ lines with $n$ points each, each class a partition of the points. Two lines of different classes meet in exactly one
point: a row and a column at their cell, a row and a symbol line because the row contains the symbol once, and two symbol lines by orthogonality. Two lines of the same class are disjoint. (A1):
fix a point $P$; it lies on $n+1$ lines, one per class, pairwise meeting only at $P$, and each containing $n-1$ further points; these $(n+1)(n-1)=n^2-1$ points are distinct, so they are all the other
points, each joined to $P$ by exactly one line. (A2): the line through $P$ in the class of $\ell$ is disjoint from $\ell$, and every line through $P$ of another class meets $\ell$. (A3): the cells
$(1,1),(1,2),(2,1)$.
\end{proof}

\begin{corollary}\label{cor:pmols}
If $q$ is a prime power, $N(q)=q-1$: for $a\in\mathbb F_q\setminus\{0\}$, the squares $L_a(x,y)=ax+y$, with rows and columns indexed by $\mathbb F_q$, are $q-1$ MOLS.
\end{corollary}
\begin{proof}
For fixed $x$, $y\mapsto ax+y$ is a bijection, and for fixed $y$, $x\mapsto ax+y$ is one since $a\ne0$. For $a\ne b$, the system $ax+y=s$, $bx+y=t$ has the unique solution $x=\frac{s-t}{a-b}$,
$y=s-ax$. With Proposition~\ref{prop:molsbound}, $N(q)=q-1$.
\end{proof}

\begin{proposition}[MacNeish, 1922]\label{prop:macneish}
$N(mn)\ge\min(N(m),N(n))$. Consequently, if $n=q_1\cdots q_t$ with $q_i$ powers of distinct primes, $N(n)\ge\min_i(q_i-1)$.
\end{proposition}
\begin{proof}
For Latin squares $L$ of order $m$ and $M$ of order $n$, define $L\otimes M$ on rows and columns $[m]\times[n]$ by $(L\otimes M)\bigl((i,i'),(j,j')\bigr)=\bigl(L(i,j),M(i',j')\bigr)$. In a row $(i,i')$, the entries
$(L(i,j),M(i',j'))$ run over all pairs as $(j,j')$ varies: it is Latin. If $L_1\perp L_2$ and $M_1\perp M_2$, then from the pair of entries of $L_1\otimes M_1$ and $L_2\otimes M_2$ at a cell one reads
$(L_1(i,j),L_2(i,j))$, which determines $(i,j)$, and $(M_1(i',j'),M_2(i',j'))$, which determines $(i',j')$: so $L_1\otimes M_1\perp L_2\otimes M_2$. Taking $\min(N(m),N(n))$ squares on each side
and pairing them gives the first claim, and Corollary~\ref{cor:pmols} the second.
\end{proof}

For $n\equiv2\pmod4$, MacNeish's bound is $N(n)\ge1$, and Euler conjectured in 1782 that there it is the truth: no two orthogonal Latin squares of order $n\equiv2\pmod 4$. The case $n=6$ is his
problem of thirty-six officers, of six ranks from six regiments, to be arranged in a square with each rank and each regiment once in every row and column; Gaston Tarry checked by hand in
1900 that it is impossible. But in 1959--1960 Bose, Shrikhande and Parker, nicknamed ``Euler's spoilers'' on the front page of the \emph{New York Times}, showed that orthogonal pairs exist
for every $n$ except $2$ and $6$. MacNeish had also conjectured that his bound is exact; it fails already at $n=10$.

\section{The Bruck--Ryser theorem}

Order $6$ is excluded twice over: no two orthogonal Latin squares, hence no projective plane. A deeper reason excludes infinitely many orders.

\begin{lemma}\label{lem:fourmatrix}
Let $n=a^2+b^2+c^2+d^2$ with integers $a,b,c,d$. The matrix
\[
A=\begin{pmatrix}a&b&c&d\\-b&a&-d&c\\-c&d&a&-b\\-d&-c&b&a\end{pmatrix}
\]
satisfies $A^{\mathsf T}A=nI_4$. So if $y=Ax$, then $y_1^2+y_2^2+y_3^2+y_4^2=n(x_1^2+x_2^2+x_3^2+x_4^2)$, and for $n\ge1$, $x=\frac1nA^{\mathsf T}y$.
\end{lemma}
\begin{proof}
The columns of $A$ each have squared length $a^2+b^2+c^2+d^2=n$, and any two of them are orthogonal: for instance, the first two give $ab-ba+cd-dc=0$, and the others are checked the same way. So
$A^{\mathsf T}A=nI$, $|Ax|^2=x^{\mathsf T}A^{\mathsf T}Ax=n|x|^2$, and $A^{-1}=\frac1nA^{\mathsf T}$. (This is Euler's four-square identity of Tome~1, Theorem~7.4, in matrix form.)
\end{proof}

\begin{theorem}[Bruck and Ryser, 1949]\label{thm:bruckryser}
If a projective plane of order $n\equiv1$ or $2\pmod4$ exists, then $n$ is a sum of two squares.
\end{theorem}
\begin{proof}
Let $v=n^2+n+1$ and let $N$ be the $v\times v$ incidence matrix, $N_{P\ell}=1$ when the point $P$ is on the line $\ell$. Row $P$ of $NN^{\mathsf T}$ counts the lines through $P$ and $P'$: it is $n+1$ for $P'=P$
and $1$ otherwise. So $NN^{\mathsf T}=nI+J$, and for $x\in\Q^v$, with $z=N^{\mathsf T}x$ and $w=x_1+\dots+x_v$,
\[
z_1^2+\dots+z_v^2=x^{\mathsf T}NN^{\mathsf T}x=n(x_1^2+\dots+x_v^2)+w^2.
\]
For $n\equiv1,2\pmod 4$, $n^2+n+1\equiv3\pmod4$, so $v+1$ is a multiple of $4$. Add a variable $x_{v+1}$ and the term $nx_{v+1}^2$ to both sides:
\[
z_1^2+\dots+z_v^2+nx_{v+1}^2=n(x_1^2+\dots+x_{v+1}^2)+w^2.\tag{$*$}
\]
By Lagrange (Tome~1, Theorem~7.4), $n=a^2+b^2+c^2+d^2$. Apply Lemma~\ref{lem:fourmatrix} to each block of four consecutive variables $(x_1,\dots,x_4),\dots,(x_{v-2},\dots,x_{v+1})$: this defines new variables
$y_1,\dots,y_{v+1}$, related to the $x_i$ by an invertible rational linear map, with $n(x_1^2+\dots+x_{v+1}^2)=y_1^2+\dots+y_{v+1}^2$. Now regard ($*$) as an identity between quadratic forms in the
independent rational variables $y_1,\dots,y_{v+1}$; every $x_i$, every $z_i$ and $w$ is a rational linear form in the $y_j$:
\[
z_1^2+\dots+z_v^2+nx_{v+1}^2=y_1^2+\dots+y_{v+1}^2+w^2.
\]
Write $z_1=c_1y_1+(\text{terms in }y_2,\dots,y_{v+1})$. If $c_1\ne1$, impose the linear relation $y_1=\frac1{1-c_1}(\text{terms in }y_2,\dots)$, which makes $z_1=y_1$; if $c_1=1$, impose $y_1=-\frac12(\text{terms in }y_2,\dots)$,
which makes $z_1=-y_1$. Either way $z_1^2=y_1^2$, and after substituting for $y_1$ we have an identity of forms in $y_2,\dots,y_{v+1}$ with $z_1^2$ and $y_1^2$ cancelled. Repeat with $z_2$ and $y_2$, and so on up to
$z_v$ and $y_v$. What remains is an identity of forms in the single variable $y_{v+1}$:
\[
nx_{v+1}^2=y_{v+1}^2+w^2,
\]
with $x_{v+1}$ and $w$ rational multiples of $y_{v+1}$. Take $y_{v+1}=1$: $nX^2=1+W^2$ with $X,W$ rational, and $X\ne0$. Clearing denominators, $n\,a^2=b^2+c^2$ with integers $a\ne0$, $b\ne0$. If a prime $p\equiv3\pmod4$
divided $n$ to an odd power, it would divide $na^2$ to an odd power, but $b^2+c^2$ to an even power (Tome~1, Theorem~7.3). So every such prime divides $n$ to an even power, and $n$ is a sum of two
squares by the same theorem.
\end{proof}

\begin{corollary}\label{cor:no6}
There is no projective plane of order $6$, $14$, $21$ or $22$.
\end{corollary}
\begin{proof}
$6\equiv14\equiv22\equiv2$ and $21\equiv1\pmod4$, and $6=2\cdot3$, $14=2\cdot7$, $21=3\cdot7$, $22=2\cdot11$ each contain a prime $\equiv3\pmod4$ to the first power, so none is a sum of two squares (Tome~1,
Theorem~7.3).
\end{proof}

The first order the theorem does not decide is $10$. In 1989 Clement Lam, Larry Thiel and Stanley Swiercz announced, after several years of computation on a Cray supercomputer, that there is no
plane of order $10$; in 2021 Curtis Bright and his colleagues confirmed it with a SAT solver producing a checkable certificate, the method of Chapter~1. Order $12$ is open: no plane is known, and no proof
of non-existence. The theorem is a statement about rational quadratic forms, and its natural home is the Hasse--Minkowski theory of Tome~12, Chapter~7; Chowla and Ryser extended it to all symmetric designs in 1950.

\section{Pairwise testing}

A program has $k$ parameters, each with $v$ possible values. Testing all $v^k$ combinations is hopeless; but experience shows that most bugs are triggered by the interaction of one or two
parameters. So we ask for a set of tests in which every pair of parameters takes every pair of values at least once.

\begin{definition}\label{def:ca}
A \emph{covering array} $\mathrm{CA}(N;k,v)$ is an $N\times k$ array with entries in a set of $v$ symbols such that for any two columns, each of the $v^2$ pairs of symbols appears in at least one row.
\end{definition}

Clearly $N\ge v^2$. Orthogonal Latin squares give the best possible arrays when there are not too many parameters.

\begin{proposition}\label{prop:camols}
If $L_1,\dots,L_{m}$ are MOLS of order $v$, the $v^2\times(m+2)$ array with rows $\bigl(x,y,L_1(x,y),\dots,L_m(x,y)\bigr)$, $(x,y)\in[v]^2$, is a $\mathrm{CA}(v^2;m+2,v)$ in which each pair appears exactly
once. In particular, for a prime power $q$ there is a $\mathrm{CA}(q^2;q+1,q)$.
\end{proposition}
\begin{proof}
The columns $x,y$: the pairs $(x,y)$ are all different. The columns $x$ and $L_k$: for fixed $x$, $y\mapsto L_k(x,y)$ is a bijection since row $x$ is Latin. Similarly for $y$ and $L_k$ (columns are Latin), and for
$L_k$, $L_l$ by orthogonality. In each case the $v^2$ rows give $v^2$ distinct pairs, hence all of them, each once.
\end{proof}

For many parameters, a random array is good, and the number of tests grows only like the logarithm of $k$.

\begin{theorem}\label{thm:caprob}
If $N\ge v^2\ln\bigl(\binom k2v^2\bigr)$, a $\mathrm{CA}(N;k,v)$ exists.
\end{theorem}
\begin{proof}
Fill an $N\times k$ array with independent uniform symbols. For two columns and a pair of symbols, the probability that no row shows that pair in those columns is $(1-v^{-2})^N<e^{-N/v^2}\le\bigl(\binom k2v^2\bigr)^{-1}$.
There are $\binom k2v^2$ such requirements, so the expected number of unmet ones is less than $1$, and some array meets them all (the first moment method, Tome~3, Chapter~12).
\end{proof}

For $v=2$ and $k=1000$ binary parameters, this gives $59$ tests instead of $2^{1000}$. For binary parameters the exact answer is known: the smallest $N$ with $\binom{N-1}{\lceil N/2\rceil}\ge k$
(Katona, 1973; Kleitman and Spencer, 1973), which is $14$ for $k=1000$; Exercise~\ref{ex:binaryca} gives the construction, and the optimality is a theorem in the spirit of Sperner's (Chapter~2).
The random proof can be made into an algorithm by choosing the rows one by one so as to cover at least the expected number of new pairs (Exercise~\ref{ex:caderand}); tools for pairwise
testing, such as the AETG system of Cohen, Dalal, Fredman and Patton (1997), work this way.

\begin{beauty}
A rector's puzzle about schoolgirls, a statistician's field trials, Euler's officers and a tester's checklist all ask for the same thing: a family of small sets in which every pair is treated
the same. The answers come from geometry over finite fields, where lines meet as Euclid wished, and the obstructions come from number theory: whether a projective plane of order $n$ exists depends,
through Lagrange's four squares, on whether $n$ is a sum of two. And at the frontier, the plane of order $10$ was ruled out by a computer, and the plane of order $12$ is still waiting.
\end{beauty}

\section{Exercises}

\begin{exercise}[$\star$]\label{ex:idempotent}
Show that a commutative quasigroup of even order $n$ cannot be idempotent ($x\circ x=x$ for all $x$). Show that in Lemma~\ref{lem:halfidem} each $x<n$ is a square $y\circ y$ for exactly two $y$, and each $x\ge n$ for none.
\end{exercise}

\begin{exercise}[$\star\star$]\label{ex:sts7}
Show that any two $\mathrm{STS}(7)$ are isomorphic (differ by a renaming of the points), so that every $\mathrm{STS}(7)$ is a Fano plane.
\end{exercise}

\begin{exercise}[$\star$]\label{ex:resolvable}
Show that a resolvable $\mathrm{STS}(v)$ requires $v\equiv3\pmod6$, and that the number of parallel classes is then $\frac{v-1}2$. Check it for Table~\ref{tab:kirkman}.
\end{exercise}

\begin{exercise}[$\star$]\label{ex:complement}
Show that the complements of the blocks of a $2$-$(v,k,\lambda)$ design with $k\le v-2$ form a $2$-$(v,v-k,\lambda')$ design, and compute $\lambda'$.
\end{exercise}

\begin{exercise}[$\star\star$]\label{ex:symmetric}
In a symmetric design ($b=v$), show that any two blocks meet in exactly $\lambda$ points. (Use the incidence matrix $N$, with $NN^{\mathsf T}=(r-\lambda)I+\lambda J$, $NJ=rJ$ and $JN=kJ$.)
\end{exercise}

\begin{exercise}[$\star$]\label{ex:affinea3}
Complete the proof of Proposition~\ref{prop:affine}(a) by showing that (A3) holds.
\end{exercise}

\begin{exercise}[$\star\star$]\label{ex:oval}
In a projective plane of order $n$, a set of points with no three on a line has at most $n+2$ points, and at most $n+1$ if $n$ is odd. Find a set of $n+1$ such points in $\mathrm{PG}(2,q)$.
\end{exercise}

\begin{exercise}[$\star$]\label{ex:brlist}
Which orders $n\le30$ are excluded by the Bruck--Ryser theorem? Which orders $n\le30$ are neither prime powers nor excluded?
\end{exercise}

\begin{exercise}[$\star$]\label{ex:binaryca}
Let $N\ge3$. Take $k=\binom{N-1}{\lceil N/2\rceil}$ columns of length $N$: each has a $0$ in row $1$, and its other rows form the characteristic vector of a $\lceil N/2\rceil$-element subset of $\{2,\dots,N\}$,
all subsets being used. Show that this is a $\mathrm{CA}(N;k,2)$.
\end{exercise}

\begin{exercise}[$\star\star$]\label{ex:caderand}
Show how to construct a $\mathrm{CA}(N;k,v)$ with $N=\bigl\lceil v^2\ln\bigl(\binom k2v^2\bigr)\bigr\rceil$ deterministically, in time polynomial in $N$, $k$ and $v$, by choosing the rows one at a time and the entries of each
row one at a time.
\end{exercise}

\section{Hints}

\begin{hint}{ex:idempotent}
Count, in an idempotent commutative quasigroup, how often a symbol appears off the diagonal.
\end{hint}
\begin{hint}{ex:sts7}
Fix the block $124$ and the blocks through $1$; then the rest is forced.
\end{hint}
\begin{hint}{ex:resolvable}
A parallel class partitions $v$ points into triples.
\end{hint}
\begin{hint}{ex:complement}
Count the blocks missing both points of a pair by inclusion--exclusion.
\end{hint}
\begin{hint}{ex:symmetric}
Show that $N$ is invertible and compute $N^{\mathsf T}N=N^{-1}(NN^{\mathsf T})N$.
\end{hint}
\begin{hint}{ex:affinea3}
Take the four points of (P3); at most two are on $\ell_\infty$.
\end{hint}
\begin{hint}{ex:oval}
The lines through one point of the set; then a point outside the set, as in Exercise~\ref{ex:capplane} of Chapter~11. For the example, use a conic.
\end{hint}
\begin{hint}{ex:brlist}
List the $n\equiv1,2\pmod 4$ and test sums of two squares with Tome~1, Theorem~7.3.
\end{hint}
\begin{hint}{ex:binaryca}
Row $1$ gives the pair $00$. For $11$, two subsets of size $\lceil N/2\rceil$ of an $(N-1)$-set must meet.
\end{hint}
\begin{hint}{ex:caderand}
Keep the set $U$ of uncovered requirements. A random row covers each with probability $v^{-2}$; fix its entries one by one, never letting the conditional expectation of the number covered decrease.
\end{hint}

\section{Solutions}

\begin{solution}{ex:idempotent}
\idea Off the diagonal, a symmetric Latin square shows each symbol an even number of times.
\writeup Each symbol $s$ appears $n$ times in the table. Off the diagonal, the cells $(x,y)$ and $(y,x)$ carry the same symbol, so each symbol appears an even number of times off the diagonal.
If the quasigroup were idempotent, $s$ would appear once on the diagonal, and $n-1$ times off it; with $n$ even, $n-1$ is odd, a contradiction. In Lemma~\ref{lem:halfidem} the diagonal contains each
$x<n$ twice ($x\circ x$ and $(n+x)\circ(n+x)$) and each $x\ge n$ never: the parities work out, since $2n-2$ and $2n$ are even.
\end{solution}

\begin{solution}{ex:sts7}
\idea Seven blocks, three through each point; once a few blocks are named, there is no choice left.
\writeup An $\mathrm{STS}(7)$ has $7$ blocks and each point is on $3$. Name the points so that $\{1,2,4\}$ is a block, and the other blocks through $1$ are $\{1,3,5\}$ and $\{1,6,7\}$ (they partition $\{2,\dots,7\}$
with $\{2,4\}$; rename to achieve this). The blocks through $2$ other than $124$ partition $\{1,3,5,6,7\}\setminus\{1\}=\{3,5,6,7\}$ into two pairs, neither of which is a pair of a block already chosen ($35$ and
$67$ are). So they are $\{3,6\},\{5,7\}$ or $\{3,7\},\{5,6\}$; swapping the names $6\leftrightarrow7$ (which keeps $167$) we may assume $\{2,3,6\}$, $\{2,5,7\}$. The two remaining blocks contain the pairs not yet
covered: $34$, $37$, $47$, $45$, $46$, $56$, which must form $\{3,4,7\}$ and $\{4,5,6\}$. So every $\mathrm{STS}(7)$ is, after renaming, the system $124,135,167,236,257,347,456$: the Fano plane.
\end{solution}

\begin{solution}{ex:resolvable}
\idea A parallel class needs $3\mid v$; with $v\equiv1,3\pmod6$ that leaves $3$.
\writeup A parallel class partitions the points into triples, so $3\mid v$; with Proposition~\ref{prop:stsnec}, $v\equiv3\pmod 6$. Each class has $\frac v3$ blocks, so there are $\frac{v(v-1)/6}{v/3}=\frac{v-1}2$
classes. For $v=15$: $7$ days of $5$ rows, $35$ blocks, and each girl walks with two others each day, $14$ in all: once with each other girl.
\end{solution}

\begin{solution}{ex:complement}
\idea A block's complement contains both $x$ and $y$ exactly when the block contains neither.
\writeup For points $x\ne y$, the blocks containing neither number $b-2r+\lambda$ (inclusion--exclusion: $b$ blocks, minus the $r$ through $x$ and the $r$ through $y$, plus the $\lambda$ through both). This does not
depend on the pair, so the complements form a $2$-$(v,v-k,b-2r+\lambda)$ design, provided $b-2r+\lambda\ge1$ and $v-k\ge2$. For the Fano plane, the complements are the $7$ quadrilaterals, a $2$-$(7,4,2)$ design.
And $b-2r+\lambda\ge1$ indeed: every complement has $v-k\ge2$ points, so contains some pair, and all pairs lie in the same number of complements.
\end{solution}

\begin{solution}{ex:symmetric}
\idea $N$ is square and invertible; conjugating $NN^{\mathsf T}$ by it gives $N^{\mathsf T}N$, which counts intersections of blocks.
\writeup Rows of $N$ are points and columns blocks. $NN^{\mathsf T}=(r-\lambda)I+\lambda J$ (Proposition~\ref{prop:params}), each row of $N$ sums to $r$ ($NJ=rJ$) and each column to $k$ ($JN=kJ$). As $b=v$, $bk=vr$ gives
$r=k$. The matrix $(r-\lambda)I+\lambda J$ has eigenvalues $r-\lambda>0$ (as $k<v$) and $r-\lambda+\lambda v>0$, so it is invertible, and so is the square matrix $N$. From $NJ=rJ$, $N^{-1}J=r^{-1}J$, so
\[
N^{\mathsf T}N=N^{-1}\bigl((r-\lambda)I+\lambda J\bigr)N=(r-\lambda)I+\lambda N^{-1}JN=(r-\lambda)I+\tfrac{\lambda k}rJ=(r-\lambda)I+\lambda J.
\]
The off-diagonal entries of $N^{\mathsf T}N$ are the sizes of intersections of two blocks: all equal to $\lambda$.
\end{solution}

\begin{solution}{ex:affinea3}
\idea Of four points with no three collinear, a line holds at most two.
\writeup Let $A,B,C,D$ be as in (P3). The line $\ell_\infty$ contains at most two of them, so at least two, say $A,B$, are affine points. If $C$ (or $D$) is affine, $A,B,C$ are three affine points not on a line. If $C,D\in\ell_\infty$,
let $X$ be the common point of $AC$ and $BD$; it is not on $\ell_\infty$ (otherwise $X=C$ and $C\in BD$, so $B,C,D$ collinear), and $A,B,X$ are not collinear (otherwise $X\in AB$, and $X\in AC$ gives $X=A$, so
$A\in BD$). So $A,B,X$ work.
\end{solution}

\begin{solution}{ex:oval}
\idea From a point of the set, the other points are seen along different lines; there are only $n+1$ lines. If the bound $n+2$ is reached, every line meets the set evenly, and parity speaks.
\writeup Let $S$ have no three collinear points and $P\in S$. The lines $PQ$, $Q\in S\setminus\{P\}$, are distinct, and there are $n+1$ lines through $P$: $|S|-1\le n+1$. If $|S|=n+2$, every line through every point of $S$
contains exactly two points of $S$, so every line meets $S$ in $0$ or $2$ points. For a point $R\notin S$, the $n+1$ lines through $R$ partition $S$ into parts of size $0$ or $2$: $n+2$ is even, so $n$ is even. In
$\mathrm{PG}(2,q)$, the conic $\{(1:t:t^2):t\in\mathbb F_q\}\cup\{(0:0:1)\}$ has $q+1$ points, and no three are collinear: three of the form $(1:t:t^2)$ give a Vandermonde determinant $\prod(t_j-t_i)\ne0$, and with
$(0:0:1)$ the determinant is $\pm(t_2-t_1)\ne0$. For $q$ odd it is a maximal such set (Segre proved in 1955 that it is then the only kind).
\end{solution}

\begin{solution}{ex:brlist}
\idea Only $n\equiv1,2\pmod4$ are concerned, and among them those divisible exactly once (or an odd number of times) by a prime $\equiv3\pmod4$.
\writeup The $n\le30$ with $n\equiv1,2\pmod 4$ are $2,5,6,9,10,13,14,17,18,21,22,25,26,29,30$. Among them, the sums of two squares are $2,5,9,10,13,17,18,25,26,29$, and the others, $6,14,21,22,30$, are excluded. The
orders $n\le30$ that are not prime powers are $6,10,12,14,15,18,20,21,22,24,26,28,30$; removing the excluded ones leaves $10$ (ruled out by computer), and $12,15,18,20,24,26,28$, all open.
\end{solution}

\begin{solution}{ex:binaryca}
\idea Each of the four pairs $00$, $11$, $01$, $10$ has its own reason: row $1$, intersection, and incomparability.
\writeup Let $A\ne B$ be two of the subsets, of size $s=\lceil N/2\rceil$, and look at their two columns. Row $1$ shows $00$. Since $2s\ge N>N-1$, $A\cap B\ne\varnothing$, and a row in $A\cap B$ shows $11$. Since $A\ne B$
have the same size, neither contains the other: a row in $A\setminus B$ shows $10$, a row in $B\setminus A$ shows $01$. For $k=1000$: $\binom{12}7=792<1000\le\binom{13}7=1716$, so $14$ rows suffice.
\end{solution}

\begin{solution}{ex:caderand}
\idea The probabilistic proof says a random row covers $|U|/v^2$ uncovered requirements on average; the method of conditional expectations finds a row doing at least as well, entry by entry.
\writeup Let $U$ be the set of requirements (two columns $c<c'$, two symbols $(s,s')$) not yet covered. For a partial row with entries fixed in a set $F$ of columns, the expected number of requirements of $U$ covered when
the other entries are uniform is a sum over $U$ of probabilities $1$, $0$, $v^{-1}$ or $v^{-2}$, according to how many of $c,c'$ are in $F$ and whether they agree with $(s,s')$; it is computable in time
$O(|U|)$. It is the average over the $v$ choices of the next entry, so one of them does not decrease it. Fixing the $k$ entries this way gives a row covering at least $|U|/v^2$ requirements of $U$. So
after each row $|U|$ is multiplied by at most $1-v^{-2}$, and after $N$ rows $|U|\le\binom k2v^2(1-v^{-2})^N<1$: the array is complete. The time is $O(Nkv|U|)$, polynomial.
\end{solution}

\section*{Sources and further reading}

\begin{itemize}
\item \emph{The stories.} T.~P.~Kirkman, \emph{On a problem in combinations}, Cambridge and Dublin Mathematical Journal 2 (1847), 191--204, and Query VI, \emph{The Lady's and Gentleman's Diary} (1850), 48; the history,
including Steiner's question of 1853 and Ozanam's cards, is in R.~Wilson and J.~J.~Watkins (eds.), \emph{Combinatorics: Ancient and Modern}, Oxford University Press, 2013 (chapters on designs and on Latin squares);
L.~Euler, \emph{Recherches sur une nouvelle esp\`ece de quarr\'es magiques} (1782), Opera Omnia I.7; G.~Tarry, \emph{Le probl\`eme des 36 officiers}, Compte Rendu de l'Association Fran\c{c}aise pour l'Avancement des
Sciences (1900--1901); the ``Euler's spoilers'' story is told in the obituary of S.~S.~Shrikhande and in R.~C.~Bose, S.~S.~Shrikhande and E.~T.~Parker, \emph{Further results on the construction of mutually orthogonal Latin
squares and the falsity of Euler's conjecture}, Canadian Journal of Mathematics 12 (1960), 189--203.
\item R.~C.~Bose, \emph{On the construction of balanced incomplete block designs}, Annals of Eugenics 9 (1939), 353--399, and \emph{On the application of the properties of Galois fields to the problem of construction of
hyper-Graeco-Latin squares}, Sankhy\=a 3 (1938), 323--338; Th.~Skolem, \emph{Some remarks on the triple systems of Steiner}, Mathematica Scandinavica 6 (1958), 273--280; D.~K.~Ray-Chaudhuri and R.~M.~Wilson, \emph{Solution of Kirkman's
schoolgirl problem}, Proceedings of Symposia in Pure Mathematics 19 (1971), 187--203; P.~Kaski and P.~R.~J.~\"Osterg\aa rd, \emph{The Steiner triple systems of order 19}, Mathematics of Computation 73 (2004), 2075--2092.
\item C.~C.~Lindner and C.~A.~Rodger, \emph{Design Theory}, 2nd ed., CRC Press, 2008: the source of our presentation of the Bose and Skolem constructions; J.~H.~van Lint and R.~M.~Wilson, \emph{A Course in Combinatorics},
2nd ed., Cambridge University Press, 2001, chapters on designs, Latin squares and projective planes, with the proof of Bruck--Ryser that we follow; P.~J.~Cameron and J.~H.~van Lint, \emph{Designs, Graphs, Codes and
their Links}, Cambridge University Press, 1991.
\item R.~H.~Bruck and H.~J.~Ryser, \emph{The nonexistence of certain finite projective planes}, Canadian Journal of Mathematics 1 (1949), 88--93; S.~Chowla and H.~J.~Ryser, \emph{Combinatorial problems}, Canadian
Journal of Mathematics 2 (1950), 93--99; C.~W.~H.~Lam, \emph{The search for a finite projective plane of order 10}, American Mathematical Monthly 98 (1991), 305--318, a beautiful account; C.~Bright, K.~K.~H.~Cheung,
B.~Stevens, I.~Kotsireas and V.~Ganesh, \emph{A SAT-based resolution of Lam's problem}, Proceedings of AAAI 2021, 3669--3676.
\item G.~O.~H.~Katona, \emph{Two applications (for search theory and truth functions) of Sperner type theorems}, Periodica Mathematica Hungarica 3 (1973), 19--26; D.~J.~Kleitman and J.~Spencer, \emph{Families of
$k$-independent sets}, Discrete Mathematics 6 (1973), 255--262; D.~M.~Cohen, S.~R.~Dalal, M.~L.~Fredman and G.~C.~Patton, \emph{The AETG system: an approach to testing based on combinatorial design}, IEEE Transactions on
Software Engineering 23 (1997), 437--444.
\end{itemize}
